\chapter{Synoptic Evolution of Tropical Cyclone Senyar}
\label{ch:synoptic_evolution}

In this chapter, we will use ERA5 data to understand the synoptic evolution of the Tropical Cyclone (TC) Senyar. This step is also important to see whether ERA5 could represent the cyclone system well before using it as the initial and boundary conditions of our experiments.

In \citet{syahreza2026brief}, synoptic evolution based on ERA5 as atmospheric reanalysis data has been discussed in detail. Similarly, we have presented different variables between three different days of the TC evolution: pre-Senyar period (24 November 2025), start period (25 November 2025), and post-Senyar period (29 November 2025). The presented variables include mean sea level pressure, specific humidty, horizontal wind speed, daily total precipitation, and sea surface temperature, as seen in Figure \ref{fig:synoptic_evol}. The shown variables are in daily mean statistics, except total rainfall that is accumulated over a day.

Prior to the start of the TC in Figure \ref{fig:synoptic_evol}(1--4), a shallow low-pressure system over the Malacca Strait is observed. This is true to the reported storm, starting as disturbance 95B over the same area. Heavy rainfall had been accumulated across North Sumatra and Malay Peninsula. The heavy rainfall is associated with low-level convergence of moist airflow across the area with source mainly coming from the South China Sea. 

At the start of the tropical cyclone in Figure \ref{fig:synoptic_evol}(6--8), a close vortex was formed. The low-level flow was still maintained such that heavy rainfall persisted accross the area. Additionally, low wind shear around the TC indicates a favorable environment to develop and maintain cyclone, as it preserves the storm structure \citep{rios2024review}. After landing in Sumatra, the tropical cyclone was downgraded to deep depression, before turning to tropical depression as it passed Malay Peninsula towards the South China Sea. At the stage of tropical depression in Figure \ref{fig:synoptic_evol}(9--12), the low pressure system had gone weak. With weaker low-level flow convergence, the rainfall mainly affected area around Malay Peninsula.

\begin{figure}[h]
\centering
\includegraphics[width=\linewidth]{../../figs/track/era5_track_roll_2.png}
\caption{Track comparison between ERA5 and observation from IBTrACS}
\label{fig:track_era5}.
\end{figure}

To further compare ERA5 representation of Tropical Cyclone Senyar with observation, track comparison is shown in Figure \ref{fig:track_era5}. In general, ERA5 is able to represent the big pattern of the event evolution with weaker intensities. The track starts on the same location and it deviates as the tropical cyclone hit land. ERA5 could not capture the intensification of the tropical cyclone core at 25 November 2025 18:00 UTC, characterized by the lowest observed sea level pressure. The wind intesity in ERA5 is also weaker as it always stays as tropical depression throughout the event. This finding on ERA5 underestimation on minimum sea level pressure and maximum wind speed is consistent with previous study on ERA5 performance representing TCs \citep{dulac2024assessing}.

\begin{figure}[h]
\centering
\includegraphics[width=\linewidth]{../../figs/synoptic_era5/synoptic.png}
\caption{Synoptic evolution of Tropical Cyclone Senyar from ERA5 on 24 November 2025 (1--4), 25 November 2025 (5--8), and 28 November 2025 (9--12). Panels (1, 5, 9) show sea surface temperature with mean sea level pressure as contour lines. Panels (2, 6, 10) show wind shear magnitude and speed by comparing height of 200 hPa and 850 hPa. Panels (3, 7, 11) show specific humidity with 10-m wind speed as stream plot. Panels (4, 8, 12) show daily rainfall over the area.}
\label{fig:synoptic_evol}.
\end{figure}